\pdfoutput=1
\documentclass[fleqn,10pt]{wlscirep}
\usepackage[utf8]{inputenc}
\usepackage[T1]{fontenc}
\usepackage{bm}
\usepackage{caption}
\usepackage{subcaption}
\usepackage{parskip}
\usepackage{booktabs}
\usepackage{multirow}
\usepackage{soul}
% \usepackage{natbib}
\usepackage{float}
\newcommand{\boldface}[1]{\boldsymbol{#1}}  
\newcommand{\bfa}{\boldface{a}}
\newcommand{\bfb}{\boldface{b}}
\newcommand{\bfc}{\boldface{c}}
\newcommand{\bfd}{\boldface{d}}
\newcommand{\bfe}{\boldface{e}}
\newcommand{\bff}{\boldface{f}}
\newcommand{\bfg}{\boldface{g}}
\newcommand{\bfh}{\boldface{h}}
\newcommand{\bfi}{\boldface{i}}
\newcommand{\bfj}{\boldface{j}}
\newcommand{\bfk}{\boldface{k}}
\newcommand{\bfl}{\boldface{l}}
\newcommand{\bfm}{\boldface{m}}
\newcommand{\bfn}{\boldface{n}}
\newcommand{\bfo}{\boldface{o}}
\newcommand{\bfp}{\boldface{p}}
\newcommand{\bfq}{\boldface{q}}
\newcommand{\bfr}{\boldface{r}}
\newcommand{\bfs}{\boldface{s}}
\newcommand{\bft}{\boldface{t}}
\newcommand{\bfu}{\boldface{u}}
\newcommand{\bfv}{\boldface{v}}
\newcommand{\bfw}{\boldface{w}}
\newcommand{\bfx}{\boldface{x}}
\newcommand{\bfy}{\boldface{y}}
\newcommand{\bfz}{\boldface{z}}
\newcommand{\bfA}{\boldface{A}}
\newcommand{\bfB}{\boldface{B}}
\newcommand{\bfC}{\boldface{C}}
\newcommand{\bfD}{\boldface{D}}
\newcommand{\bfE}{\boldface{E}}
\newcommand{\bfF}{\boldface{F}}
\newcommand{\bfG}{\boldface{G}}
\newcommand{\bfH}{\boldface{H}}
\newcommand{\bfI}{\boldface{I}}
\newcommand{\bfJ}{\boldface{J}}
\newcommand{\bfK}{\boldface{K}}
\newcommand{\bfL}{\boldface{L}}
\newcommand{\bfM}{\boldface{M}}
\newcommand{\bfN}{\boldface{N}}
\newcommand{\bfO}{\boldface{O}}
\newcommand{\bfP}{\boldface{P}}
\newcommand{\bfQ}{\boldface{Q}}
\newcommand{\bfR}{\boldface{R}}
\newcommand{\bfS}{\boldface{S}}
\newcommand{\bfT}{\boldface{T}}
\newcommand{\bfU}{\boldface{U}}
\newcommand{\bfV}{\boldface{V}}
\newcommand{\bfW}{\boldface{W}}
\newcommand{\bfX}{\boldface{X}}
\newcommand{\bfY}{\boldface{Y}}
\newcommand{\bfZ}{\boldface{Z}}
\newcommand{\bfFe}{\boldface{F}_{\text{e}}}
\newcommand{\bfCe}{\boldface{C}_{\text{e}}}
\newcommand{\bfFp}{\boldface{F}_{\text{p}}}
\newcommand{\bfFtw}{\boldface{F}_{\text{tw}}}
\newcommand{\bfFin}{\boldface{F}_{\text{in}}}
\newcommand{\bfFd}{\boldface{F}_{\text{d}}}
\newcommand{\bfepse}{\boldsymbol{\varepsilon}_{\text{e}}}
\newcommand{\bfepsp}{\boldsymbol{\varepsilon}_{\text{p}}}
\newcommand{\bfeps}{\boldsymbol{\varepsilon}}
\newcommand{\bfepsvar}{\boldsymbol{\epsilon}}
\newcommand{\bfalpha}{\boldsymbol{\alpha}}
\newcommand{\bfbeta}{\boldsymbol{\beta}}
\newcommand{\bfgamma}{\boldsymbol{\gamma}}
\newcommand{\bfdelta}{\boldsymbol{\delta}}
\newcommand{\bfepsilon}{\boldsymbol{\varepsilon}}
\newcommand{\bfzeta}{\boldsymbol{\zeta}}
\newcommand{\bfeta}{\boldsymbol{\eta}}
\newcommand{\bftheta}{\boldsymbol{\theta}}
\newcommand{\bfkappa}{\boldsymbol{\kappa}}
\newcommand{\bflambda}{\boldsymbol{\lambda}}
\newcommand{\bfrho}{\boldsymbol{\rho}}
\newcommand{\bfmu}{\boldsymbol{\mu}}
\newcommand{\bfnu}{\boldsymbol{\nu}}
\newcommand{\bfpi}{\boldsymbol{\pi}}
\newcommand{\bfxi}{\boldsymbol{\xi}}
\newcommand{\bfsigma}{\boldsymbol{\sigma}}
\newcommand{\bftau}{\boldsymbol{\tau}}
\newcommand{\bfphi}{\boldsymbol{\phi}}
\newcommand{\bfvarphi}{\boldsymbol{\varphi}}
\newcommand{\bfVarphi}{\boldsymbol{\Phi}}
\newcommand{\bfchi}{\boldsymbol{\chi}}
\newcommand{\bfomega}{\boldsymbol{\omega}}
\newcommand{\bfnull}{\boldsymbol{0}}
\newcommand{\bfGamma}{\boldsymbol{\Gamma}}
\newcommand{\bfDelta}{\boldsymbol{\Delta}}
\newcommand{\bfEta}{\boldsymbol{\Eta}}
\newcommand{\bfTheta}{\boldsymbol{\Theta}}
\newcommand{\bfLambda}{\boldsymbol{\Lambda}}
\newcommand{\bfPi}{\boldsymbol{\Pi}}
\newcommand{\bfXi}{\boldsymbol{\Xi}}
\newcommand{\bfSigma}{\boldsymbol{\Sigma}}
\newcommand{\bfPhi}{\boldsymbol{\Phi}}
\newcommand{\bfPsi}{\boldsymbol{\Psi}}
\newcommand{\bfChi}{\boldsymbol{\Chi}}
\newcommand{\bfOmega}{\boldsymbol{\Omega}}
\newcommand{\bfnabla}{\boldsymbol{\nabla}}
\newcommand{\laplace}{\boldsymbol{\Delta}}
\newcommand{\calA}{\mathcal{A}}
\newcommand{\calB}{\mathcal{B}}
\newcommand{\calC}{\mathcal{C}}
\newcommand{\calD}{\mathcal{D}}
\newcommand{\calE}{\mathcal{E}}
\newcommand{\calF}{\mathcal{F}}
\newcommand{\calG}{\mathcal{G}}
\newcommand{\calH}{\mathcal{H}}
\newcommand{\calI}{\mathcal{I}}
\newcommand{\calJ}{\mathcal{J}}
\newcommand{\calK}{\mathcal{K}}
\newcommand{\calL}{\mathcal{L}}
\newcommand{\calM}{\mathcal{M}}
\newcommand{\calN}{\mathcal{N}}
\newcommand{\calO}{\mathcal{O}}
\newcommand{\calP}{\mathcal{P}}
\newcommand{\calQ}{\mathcal{Q}}
\newcommand{\calR}{\mathbb{R}}
\newcommand{\calS}{\mathcal{S}}
\newcommand{\calT}{\mathcal{T}}
\newcommand{\calU}{\mathcal{U}}
\newcommand{\calV}{\mathcal{V}}
\newcommand{\calW}{\mathcal{W}}
\newcommand{\calX}{\mathcal{X}}
\newcommand{\calY}{\mathcal{Y}}
\newcommand{\calZ}{\mathcal{Z}}
\newcommand{\dsA}{\mathbb{A}}
\newcommand{\dsB}{\mathbb{B}}
\newcommand{\dsC}{\mathbb{C}}
\newcommand{\dsD}{\mathbb{D}}
\newcommand{\dsE}{\mathbb{E}}
\newcommand{\dsF}{\mathbb{F}}
\newcommand{\dsG}{\mathbb{G}}
\newcommand{\dsH}{\mathbb{H}}
\newcommand{\dsI}{\mathbb{I}}
\newcommand{\dsJ}{\mathbb{J}}
\newcommand{\dsK}{\mathbb{K}}
\newcommand{\dsL}{\mathbb{L}}
\newcommand{\dsM}{\mathbb{M}}
\newcommand{\dsN}{\mathbb{N}}
\newcommand{\dsO}{\mathbb{O}}
\newcommand{\dsP}{\mathbb{P}}
\newcommand{\dsQ}{\mathbb{Q}}
\newcommand{\dsR}{\mathbb{R}}
\newcommand{\dsS}{\mathbb{S}}
\newcommand{\dsT}{\mathbb{T}}
\newcommand{\dsU}{\mathbb{U}}
\newcommand{\dsV}{\mathbb{V}}
\newcommand{\dsW}{\mathbb{W}}
\newcommand{\dsX}{\mathbb{X}}
\newcommand{\dsY}{\mathbb{Y}}
\newcommand{\dsZ}{\mathbb{Z}}
\DeclareMathOperator{\divv}{div}
\DeclareMathOperator{\Curl}{Curl}
\DeclareMathOperator{\Divv}{Div}
\DeclareMathOperator{\Cof}{Cof}
\DeclareMathOperator{\cof}{cof}
\DeclareMathOperator{\Ree}{Re}
\DeclareMathOperator{\Imm}{Im}
\DeclareMathOperator{\tr}{tr}
\DeclareMathOperator{\curl}{curl}
\DeclareMathOperator{\argmin}{{arg\,min}}
\DeclareMathOperator{\arginf}{{arg\,inf}}
\DeclareMathOperator{\diag}{diag}
\DeclareMathOperator{\trace}{tr}
\DeclareMathOperator{\sign}{sign}
\DeclareMathOperator{\dev}{dev}
\DeclareMathOperator{\grad}{grad}
\DeclareMathOperator{\Grad}{Grad}
\DeclareMathOperator{\sym}{sym}
\DeclareMathOperator{\sk}{skew}
\def\Felast{\bfF_\mathbf{\!e}}
\def\Fplast{\bfF_\mathbf{\!p}}
\def\Cplast{\bfC_\mathbf{\!p}}
\DeclareSymbolFont{matha}{OML}{txmi}{m}{it}% txfonts
\DeclareMathSymbol{\varv}{\mathord}{matha}{118}
\def\epselast{\bfeps_\mathbf{\!e}}
\def\epsplast{\bfeps_\mathbf{\!p}}
\def\GLin{\mbox{{\sf GL}$(d)$}}  %{\R^{d\times d}_*}% invertible matrices
\def\Lin{\R^{d\times d}}        % all d times d matrices
\def\reff#1{(\ref{#1})}
\def\red{{\mathrm{red}}}
\def\rep{{\mathrm{rep}}}
\def\cond{{\mathrm{cond}}}
\newcommand{\kopfcolor}{\color{red}}
\newcommand{\er}{\hspace*{4cm}}
\newcommand{\n}{{\rm n}}
\newcommand{\nn}{{\rm n+1}}
\newcommand{\PsiS}{\Psi_\mathrm{S}}
\newcommand{\PsiV}{\Psi_\mathrm{V}}
\newcommand{\PsiI}{\Psi_\mathrm{I}}
\newcommand{\cA}{c_\mathrm{A}}
\newcommand{\cM}{c_\mathrm{M}}
\newcommand{\calAe}{\underset{e=1}{\overset{n_e}{\mathcal{A}}}}
\newcommand{\tcr}{\tau_{\text{crit}}}
\newcommand{\be}{\begin{equation}}
	\newcommand{\ee}{\end{equation}}
\newcommand{\beq}{\begin{eqnarray}}
	\newcommand{\eeq}{\end{eqnarray}}
\newcommand{\bem}{\begin{multline}}
	\newcommand{\eem}{\end{multline}}
\newcommand{\ba}{\begin{align}}
	\newcommand{\ea}{\end{align}}
\newcommand{\bfzero}{{\fg0}}

\usepackage[figurename=Supplementary Figure]{caption}
\usepackage[tablename=Supplementary Table]{caption}
\newcommand{\fp}[2]{\frac{\partial #1}{\partial #2} }
\newcommand{\LR}{{\qquad \Leftrightarrow \qquad }}
\newcommand{\micro}{^\mathrm{m}}
\newcommand{\RUC}{^\mathrm{RUC}}
\newenvironment{enumArabic}{
	\begin{enumerate}[label=(\arabic*),topsep=0pt]\setlength{\itemsep}{2pt}\setlength{\parskip}{0pt}
	}{\end{enumerate}
}
\newenvironment{enumRoman}{
	\begin{enumerate}[label=(\roman*),topsep=0pt]\setlength{\itemsep}{2pt}\setlength{\parskip}{0pt}
	}{\end{enumerate}
}
\newcommand{\argminWithArgs}{\operatornamewithlimits{arg\ min}}

%%%%%%%%%%%%%%%%%%%%%%%%%%%%%%%%%%%%%%%%%%%%%%%%%%%%%%%%%%%%%%%%%%%%%%%%%
%%%%%%%%%%%%%%%%%%%%%%%%%%%%%%%%%%%%%%%%%%%%%%%%%%%%%%%%%%%%%%%%%%%%%%%%%

\begin{document}
% \beginsupplement
{\sffamily\Large \textbf{Unifying the design space of truss metamaterials by generative modeling}}\\[0.2cm]
Li Zheng, Siddhant Kumar \& Dennis M. Kochmann

\vskip 0.5cm
{\sffamily\Large \textbf{\textit{Supplementary Information}}}
\flushbottom
\section{Dataset generation}
\label{sec:SI-dataGeneration}
Supplementary Figure \ref{fig:SI_dataGeneration} shows the details of the generation of a diverse truss lattice dataset. We define the truss graph in the octant within a domain $\Omega_0 = [0,1]^3 \subset \mathbb{R}^3$, which is populated into truss structures in the domain $\Omega = [-1,1]^3 \subset \mathbb{R}^3$ through reflections about the three mutually orthogonal symmetry planes. We aim to construct a sufficiently large truss design space, which exploits the wide range of achievable material properties. Starting from the five elementary truss lattices shown in Step 1 of Supplementary Fig.~\ref{fig:SI_dataGeneration}, new structures are created by randomly perturbing both the node positions and connectivities for several iterations. Supplementary Figure~\ref{fig:SI_trussAx} shows the node positions and connectivities of the five considered elementary truss structures. Node positions are altered by offsets defined in the natural coordinate system \cite{panetta2015elastic} and sampled from the uniform distribution $\lambda \sim \mathcal{U}(-0.5,0.5)$. Based on the initial truss structures, new connectivities may be introduced by removing available nodes or inserting new nodes with the following constraints: (1)~the established structure is a single connected component; (2)~all beams are shorter than $r_{\text{max}} = \sqrt{3}/2$ (times the unit side length of the unit cell); the maximum permissible length of a beam connection $r_{\text{max}}$ is chosen according to the longest connection in the initial five elementary trusses; (3)~no dangling connections exist within the structure (every node has at least two connected beams). Each geometry is perturbed for 10 iterations to generate a library that contains a wide range of truss structures with several unique topologies. From the established set, we randomly sample two lattices with repetitions allowed, which are then superimposed according to their matching nodes to yield a more diverse dataset. The full dataset considered for training the generative models contains $965{,}736$ lattices and their homogenized effective stiffness properties. 

While we did not consider specific manufacturability constraints in this work, it is possible to enforce such constraints to narrow down the design space, depending on the requirements of specific applications. For instance, geometric constraints can be imposed to limit the maximum size or number of edges of the truss graph, while printability constraints \cite{panetta2015elastic} can be considered by ensuring that every node has at least one supporting node. Such considerations can easily be incorporated into our graph-based parameterization through conditions on the structure's geometry, utilizing the graph representation. Moreover, our framework can be extended to consider structural instability or failure, e.g., by integrating buckling strength \cite{maurizi2022inverse} or compressive loading curves \cite{wang2022inverse, deng2022inverse} among the target properties towards the inverse design of buckling-resistant structures. Our generative modeling framework based on a graph representation of truss structures can be adapted to incorporate specific design considerations such as manufacturability requirements, making it a flexible framework that can be tailored to meet specific application goals. In addition, one advantage of our approach is that -- when identifying an optimal structure -- the latent space provides many structures with similar properties (nearby points in the latent space), so that secondary targets can also be implemented by choosing from a selection of potential trusses with different architectures yet similar effective stiffness.

\begin{figure}[!bp]
\centering
\includegraphics[width = \textwidth]{data-generation.png}
\caption{\textbf{Overview of the truss dataset generation.}}
\label{fig:SI_dataGeneration}
\end{figure}

\begin{figure}[ht]
\centering
\includegraphics[width = \textwidth]{ax-initial-truss.png}
\caption{\textbf{Node positions $\boldsymbol{x}$ and adjacency matrix $\boldsymbol{A}$ of each of the five considered elementary truss unit cells.}}
\label{fig:SI_trussAx}
\end{figure}

\section{Computational homogenization details}
\label{sec:SI-homogenization}

The effective stiffness tensor $\mathbb{C}$ of all truss structures in our library are computed by finite element (FE) homogenization with periodic boundary conditions \cite{glaesener2020continuum}, using on an in-house C++ FE code (available at \href{http://ae108.ethz.ch}{http://ae108.ethz.ch}). Each strut in the truss unit cell is modeled as a linear elastic Timoshenko beam with a circular cross-section. The strut radius of each unit cell is scaled to maintain a constant relative density of $\rho = 0.15$. We visualize the 3D anisotropic stiffness of truss lattices as elastic surfaces, which indicate the effective directional Young’s modulus $E(\boldsymbol{d})$ for all directions $\boldsymbol{d}\in S^2$ as
\begin{equation}
    E(\boldsymbol{d}) = \left(\sum_{i,j,k,l=1}^{3}\mathbb{C}^{-1}_{ijkl}d_i d_j d_k d_l\right)^{-1}.
\end{equation}
In the main article, we presented representative samples of novel truss lattices generated by interpolating between known structures that exhibit extreme universal anisotropy values $A^U$ \cite{ranganathan2008universal}. $A^U$ can be interpreted as a generalization of the Zener index \cite{zener1949elasticity}, which applies to structures with cubic symmetry and can be expressed as
\begin{equation}
    A^U = 5\dfrac{G_V}{G_R}+\dfrac{K_V}{K_R}-6,
\end{equation}
where $G_V$ and $K_V$ are the Voigt estimates for, respectively, the shear and bulk moduli \cite{voigt1910lehrbuch}, and $G_R$ and $K_R$ are the Reuss estimates for, respectively, the shear and bulk moduli \cite{reuss1929berechnung}. Specifically, Reuss proposed the following relations for the bulk modulus $K$ and shear modulus $G$ in terms of the compliance components $\mathbb{S}_{ijkl}$:
\begin{align}
K_R^{-1} &= (\mathbb{S}_{1111}+\mathbb{S}_{2222}+\mathbb{S}_{3333})+2(\mathbb{S}_{1122}+\mathbb{S}_{1133}+\mathbb{S}_{2233}), \\
15 G_R^{-1} &= 4(\mathbb{S}_{1111}+\mathbb{S}_{2222}+\mathbb{S}_{3333}) - 4(\mathbb{S}_{1122}+\mathbb{S}_{1133}+\mathbb{S}_{2233})+3(\mathbb{S}_{4444}+\mathbb{S}_{5555}+\mathbb{S}_{6666})   . 
\end{align}

Analogously, we adopt the Voigt average bulk and shear moduli, calculated from the anisotropic stiffness components $\mathbb{C}_{ijkl}$ as, respectively,
\begin{align}
9 K_V &= (\mathbb{C}_{1111}+\mathbb{C}_{2222}+\mathbb{C}_{3333})+2(\mathbb{C}_{1122}+\mathbb{C}_{1133}+\mathbb{C}_{2233}), \\
15 G_V &= (\mathbb{C}_{1111}+\mathbb{C}_{2222}+\mathbb{C}_{3333}) - (\mathbb{C}_{1122}+\mathbb{C}_{1133}+\mathbb{C}_{2233})+3(\mathbb{C}_{4444}+\mathbb{C}_{5555}+\mathbb{C}_{6666})    .
\end{align}

\section{Machine learning framework}
\subsection{Protocols for NN training}
\label{subsec:SI-NNprotocols}
Details of the optimized dimensions and hyperparameters (e.g., the number of hidden layers and nodes in each layer, activation functions, learning rates, etc) of the VAE model and the property predictor are provided in Supplementary Table \ref{tab:SI_NNstructure}. 
Thresholding is applied at the end of the connectivity decoder by a sigmoid function to achieve a binary connectivity matrix.
$1\%$ of the generated dataset is used for the tuning and optimization of hyperparameters. We used the PyTorch \cite{paszke2019pytorch} package throughout the implementation of the proposed generative learning framework and leveraged its automatic differentiation engine, autograd, to automatically obtain the gradients of the homogenized effective properties with respect to the structural and geometrical features towards the optimization and inverse design  of truss lattices. Distinct from deterministic ML models, the training of the VAE model often involves a trade-off between the reconstruction accuracy and the model generative ability. For example, if the KL-divergence term is too small, the model completely ignores the latent representation and fails to learn an informative latent space (known as the KL-vanishing problem \cite{bowman2015generating}). To remedy this issue, we schedule the weight of the KL-divergence term $\beta$ via the cyclical schedule \cite{fu2019cyclical} shown in Supplementary Fig. \ref{fig:SI_klAnneal}. The training process is split into several cycles, and in each cycle $\beta$ is gradually increased from 0 to 1, using a linear function after 50 epochs. For a detailed performance comparison of various annealing schemes for $\beta$, see refs.\cite{bowman2015generating, fu2019cyclical,shao2020controlvae}. 
\begin{figure}[htb]
\centering
\includegraphics[width = 0.5\textwidth]{kl-anneal.png}
\caption{\textbf{Annealing the KL-divergence term weight $\beta$ with a cyclical schedule.}}
\label{fig:SI_klAnneal}
\end{figure}
\vspace{-1.5em}

\begin{table}[]
\begin{tabular}{@{}lccccc@{}}
\toprule
\multirow{2}{*}{}& \multicolumn{2}{c}{\small \textbf{Encoder $\mathcal{Q}_{\phi}(G)$}} & \multicolumn{2}{c}{\small \textbf{Decoder $\mathcal{P}_{\theta}(G)$}}  & \multirow{2}{*}{ \small \textbf{Property predictor $\mathcal{F}_{\omega}$}}   \\ \cmidrule(lr){2-3} \cmidrule(lr){4-5}
& \footnotesize Connectivity $\boldsymbol{A}$&\footnotesize  Node positions $\boldsymbol{x}$ & \footnotesize Connectivity $\boldsymbol{A}$ & \footnotesize Node positions $\boldsymbol{x}$ &   \\ \midrule
\footnotesize Input dimensions        & \footnotesize 278  & \footnotesize 27 & \footnotesize $d_{\text{A}}(40)$ & \footnotesize $d_{\text{x}}(40)$ & \footnotesize $d(48) = d_{\text{A}}(40)+d_{\text{x}}(40)-d_{\text{Ax}}(32)$ \\ 
\footnotesize Hidden dimensions & \footnotesize $512, 512, 512, 128$ & \footnotesize $640,640,640,512$ & \footnotesize $128,512,512,512$  &\footnotesize  $256,512,640,640$  & \footnotesize $400,800,1000,400,400,200$  \\ 
\footnotesize Output dimensions       & \multicolumn{2}{c}{\footnotesize$ d(48) = d_{\text{A}}(40)+d_{\text{x}}(40)-d_{\text{Ax}}(32)$}& \footnotesize 278 & \footnotesize 27 & \footnotesize 9 \\ 
\footnotesize Activation functions & \multicolumn{2}{c}{\footnotesize ReLU} & \multicolumn{2}{c}{\footnotesize ReLU} & \footnotesize ReLU \\
\footnotesize Optimization algorithm  & \multicolumn{2}{c}{\footnotesize Adam \cite{kingma2014adam}} & \multicolumn{2}{c}{\footnotesize Adam \cite{kingma2014adam}} & \footnotesize Adam \cite{kingma2014adam} \\ 
\footnotesize Learning rate           & \multicolumn{2}{c}{\footnotesize $5 \times 10^{-4}$} & \multicolumn{2}{c}{\footnotesize $5 \times 10^{-4}$}  & \footnotesize $5 \times 10^{-4}$ \\ 
\footnotesize Batch size              & \multicolumn{2}{c}{\footnotesize 512}   & \multicolumn{2}{c}{\footnotesize 512}  & \footnotesize 512 \\ 
\footnotesize Drop out                & \multicolumn{2}{c}{\footnotesize none}  & \multicolumn{2}{c}{\footnotesize none} & \footnotesize none\\ 
\bottomrule
\end{tabular}
\caption{ \textbf{Dimensions and training hyperparameters of the optimized Variational Autoencoder and the property predictor models.}}
\label{tab:SI_NNstructure}
\end{table}

\subsection{Overlapping embedding model}
\label{subsec:overlapping-embedding-model}

In order to accommodate the interdependencies between the structural topology (represented by the adjacency matrix) and node placements (represented by node features), we modify the VAE structure by restricting which part of a latent embedding is used for each task, as illustrated in Supplementary Fig.~\ref{fig:overlapping-embedding-model}. Let $d_{\text{A}} - d_{\text{Ax}}$ be the number of latent dimensions allocated exclusively to topology, $d_{\text{x}} - d_{\text{Ax}}$ the number of dimensions exclusively allocated to geometry, and $d_{\text{Ax}}$ the number of dimensions allocated to both. To keep the overall number of trainable parameters fixed, we reduce the outputs of the encoder ($\boldsymbol{\mu}^{\text{A}}, \log\boldsymbol{\sigma}^{\text{A}} \in \mathbb{R}^{d_{\text{A}}} $, and $\boldsymbol{\mu}^{\text{x}}, \log\boldsymbol{\sigma}^{\text{x}} \in \mathbb{R}^{d_{\text{x}}}$) to obtain the final latent representations $\boldsymbol{\mu}, \log\boldsymbol{\sigma} \in \mathbb{R}^{d} $ with $ d = d_{\text{A}} + d_{\text{x}} - d_{\text{Ax}}$ as
\begin{align}
    \boldsymbol{\mu} &= \boldsymbol{\mu}^{\text{A}}_{1:d_{\text{A}}-d_{\text{Ax}}} \oplus \dfrac{1}{2}\left(\boldsymbol{\mu}^{\text{A}}_{d_{\text{A}}-d_{\text{Ax}}+1:d_{\text{A}}} + \boldsymbol{\mu}^{\text{x}}_{1:d_{\text{Ax}}}\right) \oplus \boldsymbol{\mu}^{\text{x}}_{d_{\text{Ax+1}}:d_{\text{x}}},\\
    \log\boldsymbol{\sigma} & = \log\boldsymbol{\sigma}^{\text{A}}_{1:d_{\text{A}}-d_{\text{Ax}}} \oplus \dfrac{1}{2}\left(\log\boldsymbol{\sigma}^{\text{A}}_{d_{\text{A}}-d_{\text{Ax}}+1:d_{\text{A}}} + \log\boldsymbol{\sigma}^{\text{x}}_{1:d_{\text{Ax}}}\right) \oplus \log\boldsymbol{\sigma}^{\text{x}}_{d_{\text{Ax+1}}:d_{\text{x}}},
\end{align}
where $\oplus$ denotes vector concatenation. Our rationale is that the topological and geometrical features of truss lattices are strongly correlated and, therefore, the overlapping dimensions should learn to extract the shared information between different entities, while preserving the topology-specific and geometry-specific information within their respective dimensions of latent representations. Compared to the non-overlapping model, where shared information is stored redundantly, the overlapping model, with its ability to capture joint information of different features, allows us to adjust the importance given to their interdependencies, while maintaining a constant number of training parameters. A comprehensive evaluation of the overlapping and non-overlapping models' performance across various tasks can be found in ref.\cite{lerique2020joint}. The prediction accuracy of the model with different latent space dimensions is shown in Supplementary Table~\ref{tab:overlapping-embedding-model}. We observe that increasing the shared dimensions allocated to both beyond a certain point does not significantly improve the model's performance. Therefore, we select the latent dimensions $d_\text{A} = 8$, $d_\text{Ax} = 32$, and $d_\text{x} = 8$, which strikes a reasonable balance between capturing essential features of truss structures and maintaining computational efficiency.

By decomposing different factors of variation while preserving the explicit interdependence between factors, the model captures the underlying relationships between these entities more effectively by exploiting the decomposed representations. This decomposition not only enhances the interpretability of the model \cite{lipton2018mythos} but also improves its ability to extrapolate beyond the training distribution \cite{montero2021role}. Overall, the above characteristics of the latent space make it a meaningful and structured representation of the original discrete design space, which is advantageous for the efficient optimization and generation of novel truss structures with targeted properties.
\begin{figure}[htb]
\centering
\includegraphics[width = \textwidth]{overlap-embedding-model.png}
\caption{\textbf{Schematic of the overlapping embedding model.} $d_{\text{A}} - d_{\text{Ax}}$ is the number of latent dimensions allocated exclusively to topology; $d_{\text{x}} - d_{\text{Ax}}$ is the number of dimensions allocated exclusively to geometry; $d_{\text{Ax}}$ is the number of dimensions allocated to both.}
\label{fig:overlapping-embedding-model}
\end{figure}


\begin{table}[htb]
\centering
\begin{tabular}{@{}rccccccccccccc@{}}
\toprule
\footnotesize{$R^2$-scores for}& \footnotesize{$\boldsymbol{A}$}&  \footnotesize{($x,y,z$)   }     & \footnotesize{$\mathbb{C}_{1111}$}   & \footnotesize{$\mathbb{C}_{1122}$}  & \footnotesize{$\mathbb{C}_{1133}$}   & \footnotesize{$\mathbb{C}_{2222}$}   & \footnotesize{$\mathbb{C}_{2233}$}   & \footnotesize{$\mathbb{C}_{3333}$}   & \footnotesize{$\mathbb{C}_{2323}$}   & \footnotesize{$\mathbb{C}_{3131}$}   & \footnotesize{$\mathbb{C}_{1212}$}   \\ \midrule

\scriptsize{$d_{\text{A}} = 20$, $d_{\text{Ax}} = 8$, $d_{\text{x}} = 20$} & \footnotesize{0.990} & \footnotesize{0.995} & \footnotesize{0.985} & \footnotesize{0.957} & \footnotesize{0.974} & \footnotesize{0.976} & \footnotesize{0.967} & \footnotesize{0.982} & \footnotesize{0.970} & \footnotesize{0.974} & \footnotesize{0.968} \\

\scriptsize{$d_{\text{A}} = 16$, $d_{\text{Ax}} = 16$, $d_{\text{x}} = 16$}& \footnotesize{0.990} & \footnotesize{0.996}  & \footnotesize{0.972} & \footnotesize{0.969} & \footnotesize{0.962} & \footnotesize{0.989} & \footnotesize{0.978} & \footnotesize{0.932} & \footnotesize{\textbf{0.991}} & \footnotesize{0.971} & \footnotesize{0.955} \\

\scriptsize{$d_{\text{A}} = 8, d_{\text{Ax}} = 32, d_{\text{x}} = 8$} & \footnotesize{\textbf{0.999}}  & \footnotesize{\textbf{0.999}} & \footnotesize{\textbf{0.995}} & \footnotesize{0.987} & \footnotesize{\textbf{0.989}} & \footnotesize{\textbf{0.995}} &\footnotesize{ \textbf{0.990}} & \footnotesize{\textbf{0.995}} & \footnotesize{0.983 }& \footnotesize{\textbf{0.983}} &\footnotesize{ \textbf{0.982}} \\

\scriptsize{$d_{\text{A}} = 4$, $d_{\text{Ax}} = 40$, $d_{\text{x}} = 4$} & \footnotesize{0.989} & \footnotesize{0.994} & \footnotesize{0.965} & \footnotesize{\textbf{0.990}} & \footnotesize{0.958}& \footnotesize{0.977} & \footnotesize{0.983} & \footnotesize{0.954} & \footnotesize{0.969} & \footnotesize{0.958} & \footnotesize{0.973} \\

 \bottomrule
\end{tabular}
\caption{\textbf{Comparison of the node positions reconstruction accuracy of the VAE model and the prediction accuracy of the property predictor using different latent dimension $d_\text{\normalfont A}, d_\text{\normalfont Ax}, d_\text{\normalfont x}$.} $R^2$-scores are used as the primary metric.} 
\label{tab:overlapping-embedding-model}
\end{table}

\subsection{NN model performance}
\label{subsec:SI-NNperformance}

 Supplementary Figure \ref{fig:SI_statsLatent}a and b present the distribution and statistics of the mean and standard deviation of each dimension of the latent space. To illustrate the generative capacities of the VAE model, we randomly sampled latent vectors from the prior $d$-dimensional Gaussian distribution $\boldsymbol{z} \sim \mathcal{N}(0,\boldsymbol{I})$, where $d = d_{\text{A}} + d_{\text{x}} - d_{\text{Ax}}$. The candidate latent vectors can then be decoded into novel truss structures beyond the training dataset, as shown in detail in Supplementary Fig.~\ref{fig:SI_statsLatent}c. We evaluate the fidelity of the latent space by examining the percentage of randomly sampled points that can be decoded into valid truss structures satisfying the connectivity constraint (as described in Supplementary Note~\ref{sec:SI-dataGeneration}). Further, Supplementary Figure~\ref{fig:SI_reconAcc} lists the reconstruction accuracy of node positions along with representative examples of the reconstructed vs.~true truss structures. Supplementary Figure~\ref{fig:SI_predAcc} shows the prediction accuracy of all stiffness tensor components for the property predictor $\mathcal{F}_{\omega}$. The prediction accuracies are evaluated on a separate test dataset containing 2000 lattices and their corresponding homogenized stiffness tensor components.

\begin{figure}[htb]
\centering
\includegraphics[width = \textwidth]{latent-space-stats.png}
\caption{\textbf{Distribution and statistics of the latent space.} \textbf{(a)} Distribution of the mean of the latent space coordinates. \textbf{(b)} Distribution of the standard deviation of the latent space coordinates. \textbf{(c)} Representative examples of truss structures generated by decoding from randomly-sampled points from the latent space.}
\label{fig:SI_statsLatent}
\end{figure}

\begin{figure}[htb]
\centering
\includegraphics[width = \textwidth]{recon-accuracy.png}
\caption{\textbf{Evaluation of the reconstruction accuracy of the VAE model.} \textbf{(a)} Reconstructed vs.~true 3D components $(x,y,z)$ of the node positions in the test dataset. The reconstructed node positions were computed using the decoder. Dashed lines represent the ideal lines with zero-intercept and unit-slope; the corresponding $R^2$-scores are indicated. \textbf{(b)} Representative examples of reconstructed vs.~true truss structures generated by the VAE model.}
\label{fig:SI_reconAcc}
\end{figure}

\begin{figure}[htbp]
\centering
\includegraphics[width = \textwidth]{property-prediction-accuracy.png}
\caption{\textbf{Evaluation of the prediction accuracy of the property predictor model.} Predicted vs.~true components of the stiffness tensor~$\mathbb{C}$ in the test dataset. The predicted stiffness is computed using the property predictor $\mathcal{F}_{\omega}$. All dashed lines represent the ideal line with zero-intercept and unit-slope; the corresponding $R^2$-scores are indicated.}
\label{fig:SI_predAcc}
\end{figure}

\begin{figure}[htb]
\centering
\includegraphics[width = \textwidth]{recon-outlier.png}
\caption{\textbf{Evaluation of the reconstruction accuracy of the VAE model.} \textbf{(a)} Illustration of two selected extreme outliers from Supplementary Fig.~\ref{fig:SI_reconAcc}a, showing the truss structures and their corresponding elastic surfaces. \textbf{(b)} Reconstructed vs.~true 3D components $(x, y, z)$ of the node positions in the test dataset using $\lambda_{\text{recon}}=5$. Dashed lines represent the ideal lines with zero-intercept and unit-slope; the corresponding coefficient of determination $R^2$-scores are indicated.}
\label{fig:recon-outlier}
\end{figure}


In Supplementary Fig.~\ref{fig:recon-outlier} we further examine the two extreme outliers in Supplementary Fig.~\ref{fig:SI_reconAcc}a-y by comparing the reconstructed vs.~true structure and their corresponding elastic surfaces of the two outliers. We observe that, although there is a significant difference in the reconstructed vs.~true node positions of structure (b) in Supplementary Fig.~\ref{fig:recon-outlier}a, the difference in their effective structure and properties is minor. In addition, Supplementary Figure~\ref{fig:SI_reconAcc} shows that the reconstructed node positions correlate well with the ground truth with $R^2$-values ($\approx 0.999$) close to unity for each component. Over $99.9\%$ of all instances in the test data had errors within the interval $[-0.048,0.048]$ for component $y \in [0,1]$, indicating that the overall reconstruction performance is satisfactory. Therefore, we are confident that the VAE obtains a reliable estimation of the intrinsic underlying distribution and provides a sufficiently good reconstruction of various truss structures.

Despite near-perfect overall accuracy, there are several potential reasons for the presence of extreme outliers in the correlation plots. One possible explanation is the trade-off between loss terms in this multi-task learning problem. Specifically, the VAE model is jointly trained with the property predictor as
\be\label{eq:total-loss}
\begin{aligned}
\theta,\phi,\omega \leftarrow \argminWithArgs_{\theta,\phi,\omega} \ \ 
&  
\underbrace{
\frac{1}{N}\sum_{n=1}^N \left(\left\|\bfA^\text{(n)} - \bfA^{\text{(n)}'} \right\|^2  
\ +\  
\left\|\bfx^\text{(n)} - \bfx^{\text{(n)}'} \right\|^2 \right) 
}_\text{reconstruction loss}
 + 
\underbrace{
 \frac{1}{N}\sum_{n=1}^N \left\|\bfS^{(n)} - \calF_\omega[\bfmu^{(n)}] \right\|^2  
 }_\text{property prediction loss}\\
 &
\quad +\ 
\underbrace{
\sum_{n=1}^N D_\text{KL}\bigg(\calN\left(\left[\mu_1^{(n)},\dots,\mu_d^{(n)}\right]^\intercal,\text{diag}\left(\left[{\sigma_1^{(n)}}^2,\dots,{\sigma_d^{(n)}}^2\right]^\intercal\right)\right) \parallel \calN(\boldsymbol{0},\bfI) \bigg)
}_\text{Kullback–Leibler divergence}, \\
\end{aligned}
\ee
where $\theta,\phi,\omega$ are model parameters of the VAE. Minimizing multiple loss functions simultaneously can be challenging because, e.g., the model capacity is limited or some loss terms are in conflict such as the inherent trade-off between the reconstruction robustness and generalization (given by the Kullback-Leibler divergence) in VAEs. A common strategy is to balance the losses as weighted sums of related terms during the multi-objective optimization, whereas the choice of weights will have a strong impact on the performance of the model on different tasks. In this work, we have compared the model performance with several selections of weights and selected the best model based on the property prediction accuracy, because our focus is the downstream task of optimizing truss lattices for target properties. 

Alternatively, one can choose the best model according to other criteria depending on actual goals. For example, we modify Supplementary Equation~\ref{eq:total-loss} to include a weight $\lambda_\text{recon}>0$ for the reconstruction loss as follows:
\be\label{eq:total-loss-modified}
\begin{aligned}
\theta,\phi,\omega \leftarrow \argminWithArgs_{\theta,\phi,\omega} \ \ 
&  
\underbrace{
\lambda_\text{recon}
\frac{1}{N}\sum_{n=1}^N \left(\left\|\bfA^\text{(n)} - \bfA^{\text{(n)}'} \right\|^2  
\ +\  
\left\|\bfx^\text{(n)} - \bfx^{\text{(n)}'} \right\|^2 \right) 
}_\text{reconstruction loss}
 + 
\underbrace{
 \frac{1}{N}\sum_{n=1}^N \left\|\bfS^{(n)} - \calF_\omega[\bfmu^{(n)}] \right\|^2  
 }_\text{property prediction loss}\\
 &
\quad +\ 
\underbrace{
\sum_{n=1}^N D_\text{KL}\bigg(\calN\left(\left[\mu_1^{(n)},\dots,\mu_d^{(n)}\right]^\intercal,\text{diag}\left(\left[{\sigma_1^{(n)}}^2,\dots,{\sigma_d^{(n)}}^2\right]^\intercal\right)\right) \parallel \calN(\boldsymbol{0},\bfI) \bigg)
}_\text{Kullback–Leibler divergence}. \\
\end{aligned}
\ee
After training the model with a higher weight of the reconstruction loss term ($\lambda_{\text{recon}}=5$, as opposed to $\lambda_{\text{recon}}=1$ originally), Supplementary Figure~\ref{fig:recon-outlier}b shows that the points are more centralized around the line with unit-slope and zero-intercept with fewer outliers in the correlation plots of reconstructed vs.~true node positions (evaluated on the same test dataset used in the manuscript), which, however, can lead to a compromised or degraded property predictor accuracy due to the finite total capacity of the model, as shown in Supplementary Table~\ref{tab:modified-loss-score}.


\begin{table}[htb]
\centering
\begin{tabular}{@{}ccccccccccccc@{}}
\toprule
$R^2$-scores      & $x $    & $y$     & $z$     & $\mathbb{C}_{1111}$   & $\mathbb{C}_{1122}$  & $\mathbb{C}_{1133}$   & $\mathbb{C}_{2222}$   & $\mathbb{C}_{2233}$   & $\mathbb{C}_{3333}$   & $\mathbb{C}_{2323}$   & $\mathbb{C}_{3131}$   & $\mathbb{C}_{1212}$   \\ \midrule
$\lambda_{\text{recon}} = 1$ & 0.999 & 0.999 & 0.999 & 0.995 & 0.987 & 0.989 & 0.995 & 0.990 & 0.995 & 0.983 & 0.983 & 0.982 \\
$\lambda_{\text{recon}} = 5$ & 0.999 & 0.999 & 0.999  & 0.982 & 0.961 & 0.966 & 0.981 & 0.965 & 0.979 & 0.980 & 0.981 & 0.976 \\ \bottomrule
\end{tabular}
\caption{\textbf{Comparison of the node positions' reconstruction accuracy of the VAE model and the prediction accuracy of the property predictor, using $\lambda_{\text{\normalfont recon}}=1$ and $\lambda_{\text{\normalfont recon}}=5$.}}
\label{tab:modified-loss-score}
\end{table}

\subsection{Exploration in the latent space}

\subsubsection{Sampling in the latent space}
\label{subsubsec:SI-sample-latent-space}
To illustrate that the VAE model tends to encode lattices with similar topologies and mechanical properties in the same region of the latent space, we randomly picked a starting point $\bfz_{0}$ in the latent space and sample points in its vicinity by adding Gaussian noise to the point, according to
\begin{equation}
    \bfz = \bfz_{0} + \beta \cdot \boldsymbol{\varepsilon}, \qquad  \text{with} \quad \boldsymbol{\varepsilon} \sim \mathcal{N}(\boldsymbol{0},\boldsymbol{I}),
    \label{eq:sample-latent-space}
\end{equation}
where $\beta \in [0,1]$ is a scaling factor that determines the range of sampling. Supplementary Figure~\ref{fig:sample-latent-space}c shows representative examples of structures obtained by decoding the newly sampled points $\bfz$, along with their corresponding elastic surfaces computed by FEM homogenization. We observe that randomly sampled points in the neighborhood of the starting point exhibit similar topological features and mechanical properties. This similarity can be attributed to the property predictor that aims to learn the mapping from the latent space to properties. By incorporating physical knowledge into the VAE model, the predictor encourages the model to generate similar latent representations for trusses with similar properties for better prediction performance, as demonstrated in previous work with a VAE setup \cite{gomez2018automatic} (albeit in a completely different context). 

\begin{figure}[!ht]
\centering
\includegraphics[width = \textwidth]{sample-latent-space.png}
\caption{\textbf{Representations of sampling results in the latent space.} \textbf{(a)} Two-dimensional PCA analysis of the latent space of the VAE model jointly trained with the property predictor. Colors indicate the beam radius of the truss structure. \textbf{(b)} Illustration of points sampled in the vicinity of the starting point ($\beta = 0.05$). \textbf{(c)} Representative examples of truss structures and their corresponding elastic surface generated by decoding from randomly-sampled points from the latent space. Highlighted in red in (a) and (b) is the starting point $\bfz_0$. }
\label{fig:sample-latent-space}
\end{figure}


\subsubsection{Interpolation in the latent space}
\label{subsubsec:SI-itp-latent-space}
Figure~3 of the main article shows representative examples of interpolation in the latent space. Specifically, the truss structures presented in Figure~3a were generated by traversals along different latent axes, which involves randomly selecting a truss structure, mapping it to the latent space to obtain its latent representation, and independently modifying each of the latent variables (corresponding to topology-specific, geometry-specific, and shared dimensions in the latent space) along their negative and positive directions with a fixed step size.  When interpolating between two points whose corresponding trusses exhibit extreme mechanical properties (as illustrated in Figure~3b of the manuscript), we employ \textit{spherical linear} interpolation (slerp) rather than linear interpolation (lerp) in the latent space. 

The choice of slerp over lerp is motivated by several factors. Constructing an approximately continuous latent space of truss structures provides a significant advantage, as the vector representation of truss lattices allows for the generation of novel structures by arithmetic operations. However, performing interpolation in a high-dimensional latent space with a Gaussian prior presents challenges. First, lerp assumes a straight line between points, which ignores the underlying structure of the data distribution, as schematically shown in Supplementary Fig.~\ref{fig:slerp-lerp-schemes}a below. Additionally, lerp measures the Euclidean distance between two points, which, however, does not necessarily indicate the similarity between truss structures in high-dimensional spaces. Consequently, the linearly interpolated points may jump or traverse regions of the latent space with similar representations, leading to inconsistent or unrealistic interpolations. In contrast, slerp takes into account the spherical structure of the latent space and follows the shortest arc on a hypersphere, which reduces artifacts or unnatural transitions between points. Supplementary Figure~\ref{fig:slerp-interpolation-truss} shows that linear interpolation in our truss latent space produces less coherent geometries with abrupt changes in properties, whereas slerp interpolation (as shown in Figure~3b in the main article) exhibits smoother transitions in both geometries and properties along the interpolation path. The effectiveness of slerp has been demonstrated in the context of various generative models \cite{gomez2018automatic, white2016sampling} with uniform and Gaussian priors. Overall, slerp is generally more suitable for generating smooth and meaningful interpolations in the latent space of generative models.


\begin{figure}[htb]
\centering
\includegraphics[width = 0.6\textwidth]{slerp-lerp.png}
\caption{\textbf{Illustration of the lerp and slerp interpolations with interpolation parameter $\alpha \in \{0.25, 0.5, 0.75\}$.} \textbf{(a)} \textit{Linear} interpolation: interpolated vectors are obtained by $\text{LERP}(\boldsymbol{z}_1,\boldsymbol{z}_2;\alpha) = \alpha \cdot \boldsymbol{z}_1 + (1-\alpha)\cdot \boldsymbol{z}_2$. \textbf{(b)} \textit{Spherical} interpolation: the interpolation is based on the angle between two points, and the resulting interpolation path follows a great circle on the hypersphere.}
\label{fig:slerp-lerp-schemes}
\end{figure}

\begin{figure}[htb]
\centering
\includegraphics[width = \textwidth]{lerp-interpolation.png}
\caption{\textbf{Representative examples of linear (lerp) interpolation in the latent space.} Samples are obtained by linear interpolation between two points in the latent space, whose corresponding trusses exhibit extreme directional Young's modulus $E_{11}$. Their corresponding 3D elastic surface evolution (obtained by FE homogenization) is shown along the interpolation path.}
\label{fig:slerp-interpolation-truss}
\end{figure}

\subsection{Gradient-based optimization in the latent space}
\label{subsec:SI-opt-latent-space}
For each optimization task, we first evaluate structures in the training dataset and select the 100 closest matches in terms of the target property as the initial guesses. Then, we perform gradient-based optimization for each initial guess in parallel and identify the optimal solution followed by FE homogenization-based validation of their properties. As a representative example, Supplementary Figure~\ref{fig:sample-vicinity-opt}a illustrates truss structures obtained along the optimization path of minimizing Poisson's ratio $\nu_{21}$ (as shown in Figure~4b in the main article), which showcases the evolution of the geometries and properties. 

To further emphasize that multiple similar truss structures can be generated, we have sampled in the neighborhood of the optimal solution based on Supplementary Equation \ref{eq:sample-latent-space}. Supplementary Figure~\ref{fig:sample-vicinity-opt}b shows that novel structures generated by sampling in the neighborhood of a point in the latent space exhibit similar topological features and mechanical properties (see also Supplementary Fig.~\ref{fig:sample-latent-space}). Notably, some structures generated through sampling in the latent space even display superior target properties compared with the solution we obtained through optimization, which highlights the potential of our generative framework in not only memorizing the training data but generalizing from a comprehensive truss dataset, effectively capturing the underlying mechanical relations. This finding further reinforces the notion that our generative framework is highly structured and is capable of leveraging the encoded mechanical property information within the latent space.

\begin{figure}[htbp]
\centering
\includegraphics[width = \textwidth]{sample-opt-v21.png}
\caption{\textbf{Truss structures obtained from gradient optimization to minimize Poisson's ratio $\nu_{21}$.} \textbf{(a)} Evolution of truss geometries and their elastic stiffness properties along the optimization path. \textbf{(b)} Representative examples of truss structures obtained by sampling in the vicinity of the optimal solution in (a) with the sampling factor $\beta=0.05$ using Supplementary Equation \ref{eq:sample-latent-space}. The shown elastic stiffness surfaces were obtained by FE homogenization. }
\label{fig:sample-vicinity-opt}
\end{figure}

\subsection{Computational efficiency estimates}
To demonstrate the efficiency of our approach, Supplementary Table \ref{tab:SI_runtime} presents an overview of the computational runtime, the software and hardware resources required for different tasks. Once trained offline, the property predictor can provide accurate stiffness predictions instantly. The computational time for stiffness calculation is significantly reduced compared to FEM homogenization, which provides a convenient shortcut for the optimal design of truss lattices.

\begin{table}[ht]
\centering
\begin{tabular}[t]{lccc}
\toprule
\small Tasks & \small Software & \small Hardware & \small Runtimes\\
\midrule
\small Dataset generation& \small Python&\small CPU (30 cores)$^\dagger$& \small 4 hours\\
\small Stiffness computations (of the full dataset) using FEM & \small In-house C++ FEM code & \small CPU (30 cores)$^\dagger$&\small 20 hours\\
\small Joint training of the VAE and property predictor $\mathcal{F}_{\omega}$ & \small PyTorch in Python & \small GPU$^\S$&\small 4 hours\\
\small Stiffness prediction (of 2000 trusses) using $\mathcal{F}_{\omega}$ $^\P$& \small PyTorch in Python &\small  GPU$^\S$& \small 0.001 seconds\\
\small Optimal truss design using gradient-based optimization $^\star$ & \small PyTorch in Python & \small CPU (30 cores)$^\dagger$ & \small  5 minutes\\
\bottomrule
\end{tabular}
\caption{\textbf{Overview of the computational runtime, the software and hardware resources required for different tasks.} The reported runtimes are rough average estimates. $^\dagger$Computations were performed on the Euler \MakeUppercase{\romannumeral 4} cluster of ETH Zurich with two 18-core 2.7 GHz Intel Xeon Gold 6150 processors and 192 GB of DDR4 memory at 2666 MHz. $^\S$Computations were performed on a single Nvidia Quadro RTX 6000 24 GB GDDR6. $^\P$Runtimes for the property predictor are measured for one prediction on 2000 data samples. $^\star$Runtimes for optimization via the proposed ML framework are reported for parallel computation with 100 initial guesses (averaged over the different tasks shown in the main article).}
\label{tab:SI_runtime}
\end{table}%
%%%%%%%%%%%%%%%%%%%%%%%%%%%%%%%%%%%%%%%%%%%%%%%%%%%%%%%%%%%%%%%%%%%%%%%%%%%%%%%%%%%%%%%%%%%%%%%%%%%%%%%%%%%%%%%%%%%%%%%
%%%%%%%%%%%%%%%%%%%%%%%%%%%%%%%%%%%%%%%%%%%%%%%%%%%%%%%%%%%%%%%%%%%%%%%%%%%%%%%%%%%%%%%%%%%%%%%%%%%%%%%%%%%%%%%%%%%%%%%

\bibliography{Bib}
\newpage
\appendix
\end{document}